\documentclass[11pt]{article}
\usepackage{raeez-math-template}
\numberwithin{equation}{section}
\title{The Marked Attachment at a Toric Cusp}
\author{Raeez Lorgat}
\date{}
\hypersetup{pdfsubject={Integral cusp attachment and meridian mapping tori}}
\begin{document}
\maketitle
\begin{abstract}
The toric specialization determines an actual marked map from the cusp
boundary into its smooth neighborhood. Meridian mapping tori provide
integral generators for its kernel in every degree. We compute the local
relative homology from this map, including its degree-two and degree-three
specialization. For two specified finite fillings, the first global
attachment matrix is unimodular. The next homology group is expressed
through the inclusions of two explicit meridian tori.
\end{abstract}
\section{A degeneration with marked periods}

The normal circles of a toric divisor determine its integral specialization.
We compute them for a degeneration of complex two-tori.
The central surface retains a four-dimensional fundamental class and two
triple points. Its entire angular quotient is needed for the calculation.

Let $D$ be a complex disk with coordinate $q$, and let
$c_1,c_2,c_3,c_4$ be specified nowhere-zero holomorphic functions on $D$.
Over $D^*=D\setminus\{0\}$ consider
\begin{equation}\label{c26:eq:periods}
\mathcal F_q=(\mathbb C^*)^2/
\langle (c_1(q)q,c_2(q)),(c_3(q),c_4(q)q)\rangle.
\end{equation}
The group acts by multiplication of the two coordinates.
We shrink $D$ when necessary. The positive circles of those coordinates
are $\alpha_1,\alpha_2$. The two displayed periods, in that order, are
$\beta_1,\beta_2$.
Choose logarithms of the four units on $D$.
In additive coordinates the periods form the columns of
\[
\Omega(q)=\frac{\log q}{2\pi i}I_2+B(q),
\qquad
B(q)=\frac1{2\pi i}
\begin{pmatrix}\log c_1&\log c_3\\\log c_2&\log c_4\end{pmatrix}.
\]
For small $|q|$, the imaginary part of $\Omega$ is invertible.
Thus the four periods form a real lattice, and \eqref{c26:eq:periods} is compact.
For a counterclockwise meridian, $\log q$ increases by $2\pi i$.
Its positive monodromy is
\begin{equation}\label{c26:eq:T}
T\alpha_j=\alpha_j,\qquad T\beta_j=\beta_j+\alpha_j,
\qquad
T=\begin{pmatrix}1&0&1&0\\0&1&0&1\\0&0&1&0\\0&0&0&1\end{pmatrix}.
\end{equation}
Matrices act on columns. Products of based paths traverse the right
factor first. With this convention, positive transport is conjugation
by the positive meridian. Its mapping torus is
\begin{equation}\label{c26:eq:boundary}
B_T=(F\times[0,1])/((a,1)\sim(Ta,0)),
\qquad F=\mathbb R^4/\mathbb Z^4.
\end{equation}

We retain a specified section whose lift, after an integral period
translation, has two holomorphic unit coordinates at $q=0$.
Translation by their inverses normalizes this section to $(1,1)$.
This translation commutes with the period action and extends on each
toric chart. It does not change the periods or reverse the meridian.
For application to a geometric family, \eqref{c26:eq:periods} and this
specified section must be identified with that family's punctured
neighborhood. The existence of such an identification is an input here.
The periods and fan occur in Engel's construction
\cite[Propositions 2.7--2.9, pp.~5--7]{CuspEngel}.
We use the displayed local data as hypotheses.

\section{The smooth toric filling}

Let $L=\mathbb Z^2$ and triangulate $L\otimes\mathbb R$ by the triangles
\[
[m,m+e_1,m+e_1+e_2],\qquad
[m,m+e_1+e_2,m+e_2],\qquad m\in L.
\]
In $L\oplus\mathbb Z$, take the cones generated by the three lifted
vertices $(v,1)$ and all their faces.
Both maximal cone types have determinant one in the indicated ordering.
The associated toric threefold $\widetilde N$ is therefore smooth.
On every maximal chart, its map to the disk is
\begin{equation}\label{c26:eq:chart}
q=x_0x_1x_2.
\end{equation}
This follows because the height character has value one on each ray.
The central divisor has normal crossings on these charts.
The star of the ray $(0,0,1)$ has the six primitive rays
\[
e_1,\ e_1+e_2,\ e_2,\ -e_1,\ -e_1-e_2,\ -e_2.
\]
Its component is the smooth toric surface with that fan, namely the
blowup of $\mathbb P^2$ at its three torus fixed points.

For $k=(k_1,k_2)\in L$, the period transformation is
\begin{equation}\label{c26:eq:action}
g_k(z,w,q)=
\left(q^{k_1}c_1^{k_1}c_3^{k_2}z,
q^{k_2}c_2^{k_1}c_4^{k_2}w,q\right).
\end{equation}
Its lattice map sends $(a,b,h)$ to $(a+k_1h,b+k_2h,h)$.
It translates the fan by $k$ at height one.
Multiplication by the units in \eqref{c26:eq:action} also extends on the
charts. More explicitly, each toric character transforms into a monomial
character times an integral power of the units.
Consequently \eqref{c26:eq:action} gives a holomorphic action on
$\widetilde N$ near the central divisor.

\begin{proposition}\label{c26:prop:quotient}
For a sufficiently small disk, this action is free and properly
discontinuous, and its quotient $N_c\to D$ is proper.
The total space $N_c$ is a smooth complex threefold.
Its central fibre $Z_c$ has one normalized component, with three pairs
of toric boundary curves identified and two triple points.
\end{proposition}
\begin{proof}
At $q=0$, the components through a point correspond to a nonempty finite
set of lattice vertices of one triangle or one of its faces.
A nonzero translation cannot preserve such a set.
Thus the action is free on the central fibre.
For $q\ne0$, the real logarithms of the two periods form a matrix
$(\log|q|)I_2+O(1)$. It is invertible for small $|q|$.
A nonzero period therefore cannot fix a point of $(\mathbb C^*)^2$.

Here is also a compactness check for the quotient.
Put $u=(\log|z|,\log|w|)/\log|q|$.
Period transformations translate $u$ by the lattice with matrix
$I_2+O(1/|\log|q||)$.
Every orbit has a representative in a uniformly bounded fundamental
parallelogram. Only finitely many triangles meet these parallelograms.
When $u$ belongs to a triangle, the logarithmic absolute values of its
three chart coordinates are its barycentric coordinates times $\log|q|$.
All three coordinates consequently have absolute value at most one.
Finitely many closed chart polydisks therefore contain representatives
of every punctured fibre. The normalized central component is complete
and is covered by finitely many such toric chart polydisks as well.

For proper discontinuity, a compact subset of finitely many charts has
bounded normalized logarithmic coordinates as $q$ tends to zero.
Indeed its chart coordinates have a common upper bound, and their
product is $q$, so each logarithm divided by $\log|q|$ has a uniform bound.
If a translate of this compact set meets it, the preceding invertible
period matrix bounds $k$. On an annulus away from zero the same argument
uses unnormalized logarithms. At zero, only finitely many translated
component labels can meet the compact set.
Thus only finitely many $g_k$ can occur.
The quotient is Hausdorff and smooth by freeness and proper discontinuity.
The finite compact sets of representatives prove properness over a
smaller closed disk. Quotienting the star of a lattice vertex gives
exactly the stated identifications of its six boundary curves.
\end{proof}

All estimates hold uniformly when the four units vary over a compact
parameter interval with a common disk of holomorphy and no zeros.
This observation will keep the actual units in the comparison below.

\section{The hexagon and the angular specialization}

We first take $c_i=1$. Write the resulting central fibre as $Z$.
The angular lattice is $L$, and its character lattice is
$M=\operatorname{Hom}(L,\mathbb Z)$.
Coordinates identify both lattices with $\mathbb Z^2$.
The polygon below lies in $M\otimes\mathbb R$.
Let
\begin{equation}\label{c26:eq:hexagon}
\begin{split}
H&=\{y\in\mathbb R^2: |y_1|\le1,\ |y_2|\le1,\ |y_1+y_2|\le1\},\\
Q&=\begin{pmatrix}2&-1\\-1&2\end{pmatrix},\qquad \Gamma=Q\mathbb Z^2.
\end{split}
\end{equation}
The counterclockwise vertices of $H$ are
\[
v_0=(1,0),\ v_1=(0,1),\ v_2=(-1,1),\ v_3=(-1,0),\
v_4=(0,-1),\ v_5=(1,-1).
\]
Opposite edges are paired by translations in $\Gamma$.
The three corresponding normal lines are generated by
$e_1+e_2,e_2,e_1$.
These translates tile the plane.
Indeed, the bisector for the lattice vector $Qn$, in the metric $Q^{-1}$,
is $\langle n,y\rangle=\frac12n^tQn$.
For $n=(a,b)\in\mathbb Z^2$,
\[
\max_{y\in H}\langle n,y\rangle
 =\max(|a|,|b|,|a-b|)\le a^2-ab+b^2.
\]
For example, $a^2-ab+b^2\ge3a^2/4\ge|a|$ when $|a|\ge2$.
The cases $a=0,\pm1$ follow directly by minimizing over integral $b$.
Interchanging $a,b$, and replacing $(a,b)$ by $(a,a-b)$,
gives the other two bounds.
Equality for $n=e_1,e_2,e_1+e_2$ gives the six sides of $H$.
Thus $H$ is the closed Dirichlet cell of $\Gamma$ in this metric.
Let $A_\theta=\mathbb R^2/\mathbb Z^2$ be the angular torus.
The central fibre has the following topological description:
\begin{equation}\label{c26:eq:Z}
Z=(H\times A_\theta)/\sim.
\end{equation}
Over the interior of an edge with primitive normal $n$, collapse the
circle $\mathbb Rn/\mathbb Zn$ in $A_\theta$.
At a vertex collapse the whole angular torus.
Pair opposite edges by their translations, with the same angular
coordinate in the resulting quotient circle.

Here is a coordinate construction that also fixes the sign of the polygon.
Put $\mathcal A=H\cap M=\{0,v_0,\ldots,v_5\}$.
On the dense torus of the normalized component, use coordinates
$\zeta=(\zeta_1,\zeta_2)$ and characters $\zeta^u$ for $u\in M$.
Define the negative moment coordinate
\begin{equation}\label{c26:eq:moment}
y(\zeta)=
\frac{\sum_{u\in\mathcal A}u\,|\zeta^{-u}|^2}
     {\sum_{u\in\mathcal A}|\zeta^{-u}|^2}.
\end{equation}
The homogeneous coordinates $[\zeta^{-u}]_{u\in\mathcal A}$ extend
to the normalized toric surface.
Indeed, at a vertex $u$, two adjacent exponent differences form an
integral basis. Their monomials give its affine toric chart.
Every other exponent difference lies in the same nonnegative cone.
These six charts identify the projective closure with the six-ray surface.

For completeness, write $\rho_j=\log|\zeta_j|$.
The Hessian of $\log\sum_{u\in\mathcal A}\exp(-2\langle u,\rho\rangle)$
is four times the covariance matrix of $\mathcal A$ with positive weights.
It is positive definite, since $\mathcal A$ affinely spans $M_\mathbb R$.
For $y$ in the interior of $H$, the function
\[
\log\sum_{u\in\mathcal A}\exp(-2\langle u,\rho\rangle)
       +2\langle y,\rho\rangle
\]
is strictly convex and tends to infinity as $|\rho|\to\infty$.
Its unique critical point solves \eqref{c26:eq:moment}.
On an edge the same argument has one variable.
Thus the nonnegative real part maps bijectively onto $H$.
Compactness gives a homeomorphism.
The angular stabilizer is the primitive normal circle on each edge
and the whole torus at each vertex.
This proves the quotient description before edge identifications.

The negative sign in \eqref{c26:eq:moment} makes the divisor with ray
$n\in L$ correspond to the edge $\langle n,y\rangle=1$.
Let its endpoints be $a,b\in\mathcal A$.
On that divisor, only the homogeneous coordinates for $a,b$ survive.
With the residual character $s=\zeta^{-(b-a)}$, its moment coordinate is
\[
y=\frac{a+|s|^2b}{1+|s|^2}.
\]
The opposite edge has endpoints $a-Qn,b-Qn$.
Its residual character is the same $s$ under the toric period gluing:
$\zeta\mapsto q^{-n}\zeta$ leaves $\zeta^{-(b-a)}$ unchanged because
$\langle b-a,n\rangle=0$.
Its coordinate is therefore exactly $y-Qn$.
This verifies the three edge translations, including their endpoints.
It proves \eqref{c26:eq:Z} for the actual toric quotient.

Use coordinates $(\theta,\beta)\in A_\theta\times(\mathbb R^2/\mathbb Z^2)$
on $F$. The base coordinate $y=Q\beta$ is taken modulo $\Gamma$.
The map
\begin{equation}\label{c26:eq:p}
p:F\longrightarrow Z,
\qquad (\theta,\beta)\longmapsto[y,\theta]
\end{equation}
uses any representative $y$ in the hexagon.
It is independent of the representative by the opposite-edge identifications.
It retains the two-dimensional toric orbit over the interior of $H$.

The entire positive boundary map is more informative than \eqref{c26:eq:p}.
For $0\le v\le1$, choose $y\in H$ representing $Q\beta$ and put
\begin{equation}\label{c26:eq:pv}
p_v(\theta,\beta)=[y,\theta+vQ^{-1}y].
\end{equation}
If two representatives lie on paired edges, they differ by $Qn$,
where $n$ generates the circle collapsed on those edges.
The two angular values differ by $vn$, so they define the same point.
At paired vertices all angular values are identified.
Thus \eqref{c26:eq:pv} is continuous, including the hexagon boundary.
Furthermore
\begin{equation}\label{c26:eq:endpoints}
p_0=p,\qquad p_1=pT.
\end{equation}
Indeed $Q^{-1}y$ represents $\beta$ modulo $\mathbb Z^2$.
Consequently there is an actual map
\begin{equation}\label{c26:eq:k}
k:B_T\longrightarrow Z,\qquad [\theta,\beta,v]\longmapsto p_v(\theta,\beta).
\end{equation}
The section $\theta=\beta=0$ maps to the single interior point $[0,0]$.
In particular the specified positive meridian is null-homotopic under $k$.

\section{Comparison with the smooth toric neighborhood}

We describe the polar space used in this comparison.
On a chart \eqref{c26:eq:chart}, write
$x_j=r_je^{2\pi i\phi_j}$, where $r_j\ge0$ and
$\phi_j\in\mathbb R/\mathbb Z$. Retain $\phi_j$ even when $r_j=0$.
The map to the polar disk is
\[
(r_0,r_1,r_2;\phi_0,\phi_1,\phi_2)
\longmapsto(r_0r_1r_2,\phi_0+\phi_1+\phi_2).
\]
The toric transitions and the units induce their polar transitions.
These charts define $N_c^{\log}$ and a proper forgetful map
$\tau_c:N_c^{\log}\to N_c$.
Over $q\ne0$ the forgetful map is a homeomorphism.

We use Nakayama--Ogus \cite[Theorem 5.1]{CuspNO}.
A proper, separated, exact, relatively smooth map of log analytic
spaces induces a topological fibre bundle on their polar spaces.
The target must be fine and the source relatively coherent.
Both log structures here are fine divisorial log structures.
Fine log structures are relatively coherent, and log smoothness
implies relative smoothness in this case.
Here the characteristic maps at points with $r$ branches are
$\mathbb N\to\mathbb N^r$, $1\mapsto(1,\ldots,1)$, with $1\le r\le3$.
They are exact because their group-completed image meets
$\mathbb N^r$ in the image of $\mathbb N$.
The cokernel of the group map is free of rank $r-1$.
The associated underlying chart map is smooth, so the morphism is log smooth.
Proposition~\ref{c26:prop:quotient}
gives properness and separatedness. The theorem therefore applies.
Locally, the radial rounding can be seen without an existence argument:
write $a_j\ge0$, $\min a_j=0$, and replace them by $a_j+t$,
where $t\ge0$ uniquely solves $\prod_j(a_j+t)=|q|^2$.
The square roots are the radii, and the angles remain fixed.

\begin{theorem}\label{c26:thm:actual}
For the unit-one toric filling, the polar boundary over the central
circle is $B_T$ with the marking \eqref{c26:eq:T}.
Its forgetful map to the central fibre is exactly \eqref{c26:eq:k}.
In particular \eqref{c26:eq:p} is the actual integral specialization,
after radial transport in the polar bundle.
The neighborhood $N_1$ is homeomorphic, relative to its outer boundary,
to the mapping cylinder of $k$.
\end{theorem}
\begin{proof}
The polar space restores the normal circle along a double curve and
both angular circles at a triple point.
Write $A_\phi=\mathbb R^2/\mathbb Z^2$ for the local angular torus.
Thus the polar space over the normalized toric component is
$H\times A_\phi\times S^1_v$ before the opposite-edge identifications.
Across an edge between components indexed by $m$ and $m+n$, their
local torus coordinates are $q^{-m}(z,w)$ and $q^{-m-n}(z,w)$.
At argument $2\pi v$ of $q$, the angular coordinates differ by $-vn$.
Equation~\eqref{c26:eq:moment} makes their moment representatives
differ by $-Qn$.
Hence $\theta=\phi-vQ^{-1}y$ is unchanged across the identified edges.
At $v=1$ the endpoint identification on $\theta$ is
$\theta\mapsto\theta+\beta$.
This gives \eqref{c26:eq:boundary} and \eqref{c26:eq:T} in the full lattice.
Forgetting the restored normal angles gives
$[y,\theta+vQ^{-1}y]$, proving \eqref{c26:eq:k}.

The polar bundle over $[0,\epsilon]\times S^1$ is isomorphic to the
pullback of its restriction at radius zero.
This follows from homotopy invariance of numerable fibre bundles,
applied to the radial deformation retraction.
Choose the isomorphism $h_r$ to be the identity at $r=0$.
Its central fibres have the torus group law specified above.
Let $a(r,v)$ be the inverse image of the marked section under $h_r$.
Replace $h_r(x)$ by $h_r(x+a(r,v))$ in each torus fibre.
This gives a continuous bundle isomorphism that fixes the section.
It remains the identity at $r=0$ because $a(0,v)=0$.
The forgetful map is bijective at positive radius.
At radius zero it identifies exactly the fibres of $k$.
Compactness and the Hausdorff property then identify $N_1$ with the
mapping cylinder of $k$. Radial contraction gives the specialization
and a deformation retraction onto $Z$.
\end{proof}

We verify the full period marking of this radial comparison.
Fix the ray $q=r>0$.
The phases of the characters $z,w$ extend to the polar space,
even when their absolute values vanish or have poles.
Every such character is a Laurent monomial in a toric chart.
Its phase is the same Laurent monomial in the retained chart phases.
The period action multiplies these characters by positive real numbers.
Their phases therefore descend to a map to $A_\theta$ over the full ray.

The unquotiented polar space is also a covering of this family,
with deck group $L$ given by the period translations.
At radius zero, its positive real part is tiled by $Qm+H$, $m\in L$.
The seam relation is $y_{m+n}=y_m-Qn$.
Thus $Qm+y_m$ is a coordinate on the tiling plane,
and $g_k$ acts on it by $Qk$.
At positive radius the same covering is
$(\mathbb C^*)^2\to\mathcal F_r$.

For a loop in either end fibre, record its two phase degrees and its
period deck displacement. These give a homomorphism to
$\mathbb Z^2\oplus\mathbb Z^2$.
In the central model it takes the angular circles and the polygon
periods to the four standard basis vectors.
In the nearby period parametrization
$(z,w)=\exp(2\pi i\theta)r^\beta$, it has the same values.
Both phase degrees and deck displacement remain constant in a homotopy
through the polar family.
Consequently radial transport preserves all four marked periods.
Since it also fixes the section, it preserves the marked meridian.
This establishes the full marking, beyond the conjugacy class of $T$.

We now retain the original units rather than changing the geometric input.
Put $c_j(q,t)=\exp(t\log c_j(q))$, for $0\le t\le1$.
The same fan construction, now with parameter $t$, gives the actual
family of quotients. Proposition~\ref{c26:prop:quotient} applies uniformly.
Its total space is smooth, and each local chart has coordinates
$(x_0,x_1,x_2,t)$ with $q=x_0x_1x_2$.

Here is a concrete smooth transport in this family.
Choose a finite quotient-chart cover of the compact tube and smooth
nonnegative functions $\chi_a$ subordinate to that cover, with
$\sum_a\chi_a=1$. In chart $a$, let $V_a$ vary $t$ while holding its
three toric coordinates fixed. Set
\begin{equation}\label{c26:eq:transport}
V=\sum_a\chi_a V_a,
\qquad \frac{d}{dt}\Phi_t(x)=V(\Phi_t(x)),\qquad \Phi_0(x)=x.
\end{equation}
Near the normalized section, use charts in which the section's two
unit coordinates are fixed. Choose the cover and cutoffs accordingly.
Then $V$ follows the section there. Everywhere $dq(V)=0$ and
$dt(V)=1$. We check a stronger property than tangency to the divisor.
On an overlap, a local branch coordinate has the form $x_j=u y_k$,
where $u$ is a nowhere-zero function, holomorphic in the fibre variables
and smooth in $t$.
Holding the $y$ coordinates fixed gives
$V_a(x_j)=x_j u^{-1}\partial_tu$.
Thus each $V_a(x_j)$ is smoothly divisible by $x_j$.
Multiplication by a smooth cutoff and addition preserve this property.
Compactness ensures that \eqref{c26:eq:transport} exists for the whole interval.
Its reverse flow is its inverse. It gives smooth diffeomorphisms preserving
$q$, the central strata, the complex orientation, and the specified section.
The flow lifts to the polar spaces. In a local branch coordinate it
has the form $V(x_j)=x_j f_j$ for a smooth complex function $f_j$.
The lifted equations are
\[
\dot r_j=r_j\operatorname{Re}f_j,\qquad
\dot\phi_j=\operatorname{Im}f_j/(2\pi).
\]
They extend to $r_j=0$ and give the actual map of the normal circles.

This construction specifies a map, once the chart cutoffs are chosen.
Convex interpolation of two such fields gives an isotopy with the same
properties. On the lattice its transport is the identity in the
continuously marked periods. On the boundary fundamental group it also
retains the meridian, because it retains the section.
Thus it introduces no meridian translation.

The boundary and specialization for the given units are the conjugates
of \eqref{c26:eq:k} and \eqref{c26:eq:p} by this actual flow and radial
polar transport. A comparison with the period parametrization of
\eqref{c26:eq:periods} has the same action on the full boundary fundamental
group. The universal cover of $B_T$ is $\mathbb R^5$, with affine deck
transformations. Equivariant lifts of the two comparisons admit the
straight-line homotopy between them. It descends because affine maps
preserve convex combinations. This supplies the boundary homotopy in
the specified period marking, including its meridian.

\section{Finite cellular chains and the specialization matrix}

We now give cells on the quotient itself.
The hexagon torus $H/\Gamma$ has two vertices $A,B$, three edges
$e_0,e_1,e_2$ oriented from $A$ to $B$, and one oriented face $h$.
Here $A$ is the image of $v_0,v_2,v_4$.
The edge $e_0$ is $[v_0,v_1]$, $e_1$ is $[v_2,v_1]$, and
$e_2$ is $[v_2,v_3]$.
The face boundary has word
\[
e_0e_1^{-1}e_2e_0^{-1}e_1e_2^{-1},
\]
read along the oriented polygon; it has zero cellular boundary.
The two period cycles are
\begin{equation}\label{c26:eq:beta-cells}
\beta_1=e_1-e_0,\qquad \beta_2=e_0-e_2.
\end{equation}
Their lifts have displacements $Qe_1$ and $Qe_2$.

Give $A_\theta$ its product circle cells
$1,a_1,a_2,a_{12}$, with zero cellular differential.
Use base-first product orientations.
The actual product complex $D=C_*(F;\mathbb Z)$ has cells
$b\otimes a_I$ for each base cell $b$ and each angular cell $a_I$.
Its only boundaries are
\begin{equation}\label{c26:eq:dD}
d(e_j\otimes a_I)=(B-A)\otimes a_I.
\end{equation}
Its ranks in degrees zero through four are $(2,7,9,5,1)$.
This is a product CW decomposition of the four-torus, rather than a
replacement by its homology groups.

On each quotient edge, the angular character has the following degree row:
\begin{equation}\label{c26:eq:characters}
\lambda_0=(-1,1),\qquad \lambda_1=(-1,0),\qquad \lambda_2=(0,1).
\end{equation}
Their kernels are the indicated primitive normal circles.
These are actual cellular maps $A_\theta\to S^1$.
They send the angular vertex to the circle vertex, and the entire
two-dimensional angular cell into the target one-skeleton.
Consequently the quotient \eqref{c26:eq:Z} has cells
\[
\begin{array}{c|l}
0&A,B\\
1&e_0,e_1,e_2\\
2&E_0=e_0\otimes\eta_0,\ E_1=e_1\otimes\eta_1,\ E_2=e_2\otimes\eta_2,\ h\\
3&h\otimes a_1,\ h\otimes a_2\\
4&h\otimes a_{12}.
\end{array}
\]
These cells have explicit characteristic maps.
For $E_j$, use $[0,1]\times[0,1]$, with the second coordinate
parametrizing the residual circle and both endpoint circles collapsed.
For each $h\otimes a_I$, use the product of the hexagon's characteristic
disk with the angular characteristic cube.
Its boundary maps by the quotient characters on the six polygon edges.
Those characters are cellular on the angular product cells.
Attaching these products in increasing angular degree therefore gives
a finite CW structure on the quotient space itself.
Its chain complex $Q_*=C_*(Z;\mathbb Z)$ has ranks $(2,3,4,2,1)$ and
\begin{equation}\label{c26:eq:dQ}
d_1=\begin{pmatrix}-1&-1&-1\\1&1&1\end{pmatrix},\qquad d_n=0\quad(n\ge2).
\end{equation}
For the higher face boundaries, each pair of opposite edges has the same
quotient character and opposite boundary signs. Their contributions cancel.
The suspended edge circles have both endpoint circles collapsed.
These facts prove \eqref{c26:eq:dQ} directly from the attaching maps.

The map $p$ is cellular. Its complete cellular matrix is specified by
\begin{equation}\label{c26:eq:cellularp}
\begin{aligned}
p_*(A\otimes1)&=A,&p_*(B\otimes1)&=B,\\
p_*(e_j\otimes1)&=e_j,&
p_*(e_j\otimes a_i)&=\lambda_j(a_i)E_j,\\
p_*(h\otimes a_I)&=h\otimes a_I.&&
\end{aligned}
\end{equation}
All positive angular cells over $A,B$, and all $e_j\otimes a_{12}$,
have image zero in their cellular degree.
Equations \eqref{c26:eq:dD}--\eqref{c26:eq:cellularp} show $dp_*=p_*d$
over $\mathbb Z$, on every generator.

The base changes from $(A,B,e_0,e_1,e_2,h)$ to
$A,B-A,e_0,\beta_1,\beta_2,h$ are unimodular.
The pair $e_0\mapsto B-A$ contracts.
Tensoring this contraction with the angular circle cells gives an
integral contraction of the corresponding part of $D$.
The target has the same single contractible pair and zero differential
on its complement. Hence
\begin{equation}\label{c26:eq:homology}
H_n(Z;\mathbb Z)=\mathbb Z,\ \mathbb Z^2,\ \mathbb Z^4,\
\mathbb Z^2,\ \mathbb Z
\quad(0\le n\le4).
\end{equation}
There is no integral torsion.

Use the ordered fibre lattice $(\alpha_1,\alpha_2,\beta_1,\beta_2)$,
and lexicographic exterior bases. Use \eqref{c26:eq:beta-cells} in target
degree one, $(E_0,E_1,E_2,h)$ in degree two, and the displayed cells
in degrees three and four. Then
\begin{equation}\label{c26:eq:homologyp}
\begin{gathered}
p_{*,1}=\begin{pmatrix}0&0&1&0\\0&0&0&1\end{pmatrix},\\
p_{*,2}=\begin{pmatrix}
0&-1&1&1&-1&0\\0&1&0&0&0&0\\0&0&0&0&1&0\\0&0&0&0&0&1
\end{pmatrix},\qquad
p_{*,3}=\begin{pmatrix}0&0&1&0\\0&0&0&1\end{pmatrix},\\
p_{*,0}=1,\qquad p_{*,4}=1.
\end{gathered}
\end{equation}
For example, $\alpha_i\wedge\beta_j$ is represented by
$-\beta_j\otimes a_i$ in the base-first cells.
This fixes the signs of the middle matrix.
Each displayed matrix has a maximal minor equal to $1$ or $-1$.
All its nonzero Smith factors are therefore one.

\begin{theorem}\label{c26:thm:integral}
The specialization is surjective on integral homology in every degree.
Its kernels in positive degrees are
\[
\begin{array}{c|l}
1&\mathbb Z\alpha_1\oplus\mathbb Z\alpha_2\\
2&\mathbb Z(\alpha_1\wedge\alpha_2)\oplus
\mathbb Z(\alpha_1\wedge\beta_2-\alpha_2\wedge\beta_1)\\
3&\mathbb Z(\alpha_1\wedge\alpha_2\wedge\beta_1)\oplus
\mathbb Z(\alpha_1\wedge\alpha_2\wedge\beta_2)\\
4&0.
\end{array}
\]
In every degree its dual map has image exactly
$H^n(F;\mathbb Z)^T$. This remains true for the specified units,
under the marking supplied by \eqref{c26:eq:transport}.
\end{theorem}
\begin{proof}
The integral kernels and surjectivity follow by solving the displayed
matrices, with integer coefficients. Their kernels equal
$(\bigwedge^nT-I)H_n(F;\mathbb Z)$: expand \eqref{c26:eq:T} in each
exterior degree. In degree two, the image is generated by
$\alpha_1\wedge\alpha_2$ and
$\alpha_1\wedge\beta_2-\alpha_2\wedge\beta_1$.
The extra $\alpha_1\wedge\alpha_2$ term from $T(\beta_1\wedge\beta_2)$
does not change that integral span. The other degrees are immediate.
The kernels are saturated, since all specialization Smith factors are one.
Dualizing the integral quotient therefore gives exactly the invariant
annihilator, with no index or torsion correction.
The actual transport of the toric family preserves these chain comparisons.
\end{proof}

The ordered basis $(\alpha_1,\alpha_2,\beta_1,\beta_2)$ has the negative
of the complex orientation. The base-first top cell
$h\otimes a_{12}$ has the same negative sign relative to the toric
complex orientation. Thus $p_{*,4}=1$ has geometric degree $+1$.

\section{The boundary homotopy and an integral marking}

For normalized singular cubical chains, define an interval-first operator
\begin{equation}\label{c26:eq:H}
(H\sigma)(v,t_1,\ldots,t_n)=p_v(\sigma(t_1,\ldots,t_n)).
\end{equation}
It is a map on actual cubes, in every degree. Its boundary identity is
\begin{equation}\label{c26:eq:chainhomotopy}
dH+Hd=p_*T_*-p_*.
\end{equation}
The mapping-torus cone convention is
$d(a,b)=(da+(T_*-I)b,-db)$.
The actual boundary retraction is consequently represented by
\begin{equation}\label{c26:eq:Kchain}
K(a,b)=p_*a+Hb.
\end{equation}
Substitution of \eqref{c26:eq:chainhomotopy} proves that it is a chain map.
Cellular characteristic maps, their subdivisions, and their prisms
compare \eqref{c26:eq:cellularp} with these cubical maps.
In particular the homology matrix alone does not replace $H$.

Write $\Lambda=\bigoplus_{j=1}^4\mathbb Ze_j$, and set
\begin{equation}\label{c26:eq:commonbasis}
\alpha_1=e_1-e_3,\quad \alpha_2=e_3-e_2,\quad
\beta_1=e_3,\quad\beta_2=e_4.
\end{equation}
Their column matrix and inverse are
\[
L=\begin{pmatrix}1&0&0&0\\0&-1&0&0\\-1&1&1&0\\0&0&0&1\end{pmatrix},
\qquad L^{-1}=\begin{pmatrix}1&0&0&0\\0&-1&0&0\\1&1&1&0\\0&0&0&1\end{pmatrix}.
\]
The determinant is $-1$. Conjugation gives the actual positive cusp matrix
\begin{equation}\label{c26:eq:nativemonodromy}
LTL^{-1}=T_\infty=
\begin{pmatrix}2&1&1&0\\0&1&0&-1\\-1&-1&0&1\\0&0&0&1\end{pmatrix}.
\end{equation}
In the $e$ basis, the specialization matrices are
\begin{equation}\label{c26:eq:nativemaps}
\begin{gathered}
p_{*,1}^{e}=\begin{pmatrix}1&1&1&0\\0&0&0&1\end{pmatrix},\\
p_{*,2}^{e}=\begin{pmatrix}
0&-1&1&-1&1&0\\1&1&0&0&0&0\\0&0&0&0&-1&0\\0&0&1&0&1&1
\end{pmatrix},\qquad
p_{*,3}^{e}=\begin{pmatrix}0&1&1&0\\0&1&0&-1\end{pmatrix},\qquad
p_{*,4}^{e}=-1.
\end{gathered}
\end{equation}
The signs follow from $p_{*,n}^{e}=p_{*,n}\bigwedge^nL^{-1}$.

\begin{remark}
The two-torus $H/\Gamma$ has zero fourth homology and Euler characteristic.
The actual fibre has $H_4(Z;\mathbb Z)=\mathbb Z$ and $\chi(Z)=2$.
Thus collapsing every angular fibre changes the homotopy type
and the local attachment.
\end{remark}

\begin{thebibliography}{9}
\bibitem{CuspEngel}
P.~Engel, \emph{Complex structures on $S^6$},
\href{https://philip-engel.github.io/S6.pdf}{author manuscript},
Propositions~2.7--2.9, pp.~5--7.
\bibitem{CuspNO}
C.~Nakayama and A.~Ogus,
\emph{Relative rounding in toric and logarithmic geometry},
Geometry \& Topology \textbf{14} (2010), 2189--2241,
\href{https://doi.org/10.2140/gt.2010.14.2189}{doi:10.2140/gt.2010.14.2189},
Theorem~5.1.
\end{thebibliography}

\section{The marked map into the cusp neighborhood}

Retain the toric family, units, logarithms, section, and lattice
of Sections~1--6. Restrict to the compact tube $N_c$ over $|q|\leq\epsilon$.
Its boundary is smooth because the punctured family is a submersion.
Write $a=L(\theta,\beta)^t$ in the common $e$ marking.
For $q(v)=\epsilon e^{2\pi iv}$, use the logarithm
$\log q(v)=\log\epsilon+2\pi iv$ and set
\begin{equation}\label{ga:eq:j}
j_\infty[a,v]=
\left[\exp\bigl(2\pi i(\theta+\Omega(q(v))\beta)\bigr),q(v)\right]
\in\partial N_c.
\end{equation}
Exponentiation is coordinatewise. Brackets on the right mean the
quotient by the two specified periods.
Changing $\theta$ by an integer has no effect.
Changing $\beta$ by an integer multiplies by the corresponding periods.
At $v=1$, the matrix $\Omega$ is its value at zero plus $I_2$.
Thus \eqref{ga:eq:j} respects $[a,1]=[T_\infty a,0]$.
The same identity holds for derivatives across this seam.
The real period matrix is invertible, so this is a smooth diffeomorphism.
It carries the marked meridian to the specified section over the circle.

Give $F$ its complex orientation, which agrees with the $e$ orientation.
Give the disk its complex orientation and each boundary its outward-first
orientation. The boundary orientation of $N_c$ is fibre-first followed
by the positive meridian. The boundary of a punctured base has the
opposite meridian orientation. Thus \eqref{ga:eq:j} has the required
orientation reversal for attachment to the flat bundle over that base.

On integral normalized singular cubical chains define the actual map
\begin{equation}\label{ga:eq:kappa}
\kappa_\infty:C_*(B_{T_\infty})\longrightarrow C_*(N_c),
\qquad \sigma\longmapsto j_\infty\circ\sigma.
\end{equation}
This fixes the carrier and the map in every degree.
Finite cellular models of the compact pairs give the corresponding
finite chain map after cellular approximation and subdivision.

We relate this map to the toric cells without losing its meridian.
Use $L^{-1}$ to express the boundary in the $(\alpha,\beta)$ marking.
For $C=C_*^{\mathrm{cub}}(F;\mathbb Z)$ put
\[
R_n=C_n\oplus C_{n-1},\qquad
d_R(a,b)=(da+(T_*-I)b,-db).
\]
The interval-first prism comparison identifies $R$ with chains of $B_T$.
The toric forgetful map is represented by
\begin{equation}\label{ga:eq:K}
K:R_n\longrightarrow C_n^{\mathrm{cub}}(Z;\mathbb Z),
\qquad K(a,b)=p_*a+Hb,
\qquad (H\sigma)(v,t)=p_v(\sigma(t)).
\end{equation}
Here $p_v(\theta,\beta)=[y,\theta+vQ^{-1}y]$,
where $y$ represents $Q\beta$ in $H$.
The identity $dH+Hd=p_*T_*-p_*$ proves $dK=Kd_R$.

Let $j_r$ be the outer boundary identification from radial polar transport.
The self-map $j_r^{-1}j_\infty$ has the identity action on the full
based fundamental group, by the phase, deck, and section marking.
Choose equivariant lifts to the affine universal cover of $B_T$.
Their straight-line interpolation supplies a boundary homotopy.
Following that homotopy by the radial retraction gives a chain homotopy
\begin{equation}\label{ga:eq:comparison}
r_*\kappa_\infty-K=dS+Sd_R,
\end{equation}
with the prism comparison understood on the source.
The homotopy is fixed on the section.
It is defined by actual maps, rather than by an equality of homology ranks.
For nontrivial units, the smooth unit transport carries this construction
to the original family and its continuously specified marking.

The fibre part of \eqref{ga:eq:K} in degrees two and three is
\begin{equation}\label{ga:eq:p23}
P_2=\begin{pmatrix}
0&-1&1&1&-1&0\\0&1&0&0&0&0\\0&0&0&0&1&0\\0&0&0&0&0&1
\end{pmatrix},\qquad
P_3=\begin{pmatrix}0&0&1&0\\0&0&0&1\end{pmatrix}.
\end{equation}
Their source bases are the ordered exterior bases of
$(\alpha_1,\alpha_2,\beta_1,\beta_2)$.
Their target bases are $(E_0,E_1,E_2,h)$ and
$(h\otimes a_1,h\otimes a_2)$.
In the $e$ basis these matrices become $P_n\bigwedge^n L^{-1}$.
The term $H$ in \eqref{ga:eq:K} remains part of the boundary map.

\section{Integral meridian cycles}

Put $A_n=\bigwedge^n\mathbb Z^4$ and $N_n=\bigwedge^nT-I$.
Set $A_n=0$ outside $0\leq n\leq4$.
Write $\varepsilon_1,\varepsilon_2$ for the standard vectors in
the angular and polygon coordinate lattices $\mathbb Z^2$.
The mapping-torus chain sequence gives
\begin{equation}\label{ga:eq:wang}
0\longrightarrow A_n/N_nA_n\longrightarrow H_n(B_T;\mathbb Z)
\mathrel{\mathop{\longrightarrow}^{\pi}}\ker N_{n-1}\longrightarrow0.
\end{equation}
The left map is inclusion of the fibre.
Section~5 identifies its source with $H_n(Z;\mathbb Z)$ through $p_*$.
We now give an integral splitting for which $k_*$ vanishes on the right summand.

For a $T$-invariant oriented subtorus $V\subset F$, let $M(V)$ be its
mapping torus inside $B_T$. Orient it by the positive interval first,
followed by the stated basis of $V$.
Its fundamental class projects under $\pi$ to $[V]$.
All exterior products below are in the displayed order.
The following subtorus classes give integral bases of $\ker N_r$:
\begin{equation}\label{ga:eq:invariants}
\begin{array}{c|l}
r&[V]\\\hline
0&1\\
1&\alpha_1,\ \alpha_2\\
2&\alpha_1\alpha_2,\ \alpha_1\beta_1,\ \alpha_2\beta_2,
 (\alpha_1+\alpha_2)(\beta_1+\beta_2)\\
3&\alpha_1\alpha_2\beta_1,\ \alpha_1\alpha_2\beta_2\\
4&\alpha_1\alpha_2\beta_1\beta_2.
\end{array}
\end{equation}
Here juxtaposition denotes an exterior product.
For $r=0$, $M(V)$ is the marked meridian.
Each displayed sublattice is primitive and $T$-stable.
The restricted monodromy has determinant one.

To check the degree-two basis, expand $T\beta_j=\beta_j+\alpha_j$.
An invariant two-vector has no $\beta_1\beta_2$ term, and its coefficients
of $\alpha_1\beta_2$ and $\alpha_2\beta_1$ are equal.
The four displayed vectors form an integral basis of these solutions.
In degree three the only invariant vectors are integral combinations of
$\alpha_1\alpha_2\beta_1$ and $\alpha_1\alpha_2\beta_2$.
The remaining degrees follow directly from the same expansion.

Each class $[M(V)]$ in \eqref{ga:eq:invariants} maps to zero under $k_*$.
For purely angular $V$, the map is independent of the meridian coordinate.
For $V=\langle\alpha_j,\beta_j\rangle$, its base lies on the closed
line $\beta=t\varepsilon_j$. Representatives in $H$ have
$Q^{-1}y=(t-m)\varepsilon_j$ for an integer $m$ on each segment.
The angular shift $v(t-m)\varepsilon_j$ remains in the angular direction of $V$.
Thus its image lies in a two-dimensional quotient of rectangles times circles.
For $V=\langle\alpha_1+\alpha_2,\beta_1+\beta_2\rangle$, the same
argument uses the line $t(\varepsilon_1+\varepsilon_2)$ and seam translations
by $\varepsilon_1+\varepsilon_2$.
Its image also has a finite CW description of dimension at most two.
The three-dimensional subtori in the table have the entire angular torus
and a one-dimensional base. Their images have dimension at most three.
Finally, $k$ maps the five-dimensional class $[B_T]$ into the four-dimensional $Z$.
These factorizations prove the asserted vanishing in every degree.

\begin{theorem}\label{ga:thm:boundary}
The boundary homology is free, with ranks $(1,3,6,6,3,1)$.
For $0\leq n\leq4$, the map induced by the actual cusp attachment is
\begin{equation}\label{ga:eq:split}
H_n(B_T;\mathbb Z)=H_n(Z;\mathbb Z)\oplus\ker N_{n-1}
\mathrel{\mathop{\longrightarrow}^{k_*}}H_n(Z;\mathbb Z),
\qquad (z,u)\longmapsto z.
\end{equation}
The right summand is represented by the mapping-torus cycles above.
In degree five the map is zero.
\end{theorem}
\begin{proof}
The fibre quotient in \eqref{ga:eq:wang} maps isomorphically to $H_n(Z)$.
The mapping-torus cycles give a section of $\pi$, since their projected
classes form the integral basis \eqref{ga:eq:invariants}.
Their images under $k_*$ vanish by the explicit factorizations.
This proves \eqref{ga:eq:split}. Both summands are free.
Their ranks give the stated boundary ranks.
Equation~\eqref{ga:eq:comparison} identifies this result with
$r_*\kappa_{\infty*}$ in the original units and marking.
\end{proof}

In particular, the degree-two kernel has the actual torus basis
\begin{equation}\label{ga:eq:kernel2}
z_j=[M(\langle\alpha_j\rangle)],\qquad j=1,2.
\end{equation}
Explicitly, their fundamental squares are $(v,s)\mapsto[s\alpha_j,v]$.
Their boundaries cancel because $T\alpha_j=\alpha_j$.
The cusp retraction sends each square to $[0,s\varepsilon_j]$, independently of $v$.
The degree-three kernel has the four mapping-torus classes from row $r=2$.
These specify the meridian contributions as well as the fibre specialization.

\section{The integral local attachment cone}

For a chain map $f:B\to A$, use
$\operatorname{Cone}(f)_n=A_n\oplus B_{n-1}$ and
$d(a,b)=(da+fb,-db)$.
For the boundary inclusion \eqref{ga:eq:kappa}, this cone computes
$H_*(N_c,\partial N_c;\mathbb Z)$.
Indeed, inclusion of a boundary subcomplex is injective on cellular chains.
The cone maps onto the relative quotient, and its kernel is an identity cone.
Its contraction is $(c,b)\mapsto(0,c)$ in the identity-cone coordinates.

We describe the integral reduction used in the calculation.
A finite free chain complex with free homology splits as its homology
with zero differential plus contractible pairs.
To prove this, split the two exact sequences
\[
0\to Z_n\to C_n\to B_{n-1}\to0,
\qquad 0\to B_n\to Z_n\to H_n\to0.
\]
Here $Z_n=\ker d_n$ and $B_n=\operatorname{im}d_{n+1}$ in the complex $C$.
The quotients are free, so both sequences split over $\mathbb Z$.
The differential on the remaining pairs is the identity.
These splittings give maps $i,p,h$ with $pi=I$ and $I-ip=dh+hd$.
For a chain map $f:B\to A$, put $f_H=p_Afi_B$.
Then
\[
f-i_Af_Hp_B=dU+Ud,\qquad U=h_Af+i_Ap_Afh_B.
\]
Mapping cones of homotopic maps are isomorphic by their triangular
change of coordinates. Thus this operation preserves integral cone homology.

Both complexes here have free homology by Theorem~\ref{ga:thm:boundary}
and the local specialization theorem.
Consequently the actual local cone reduces to a finite free complex $D$.
Its degree-zero through degree-six ranks are
\begin{equation}\label{ga:eq:ranks}
(1,3,7,8,7,3,1).
\end{equation}
Write $r_n=(1,2,4,2,1)$ for the target ranks and
$b_n=(1,3,6,6,3,1)$ for the boundary ranks, with zero outside these ranges.
In the ordered coordinates $D_n=\mathbb Z^{r_n}\oplus\mathbb Z^{b_{n-1}}$,
the only nonzero block of $d_n$ is the upper-right block
\begin{equation}\label{ga:eq:reduced-d}
\begin{pmatrix}I_{r_{n-1}}&0\end{pmatrix}
:\mathbb Z^{b_{n-1}}\longrightarrow\mathbb Z^{r_{n-1}},
\qquad 1\leq n\leq5.
\end{equation}
The differential $d_6$ vanishes.
Thus the nonzero Smith factors of $d_1,\ldots,d_5$ are respectively
$1$, $(1,1)$, $(1,1,1,1)$, $(1,1)$, and $1$.

\begin{theorem}\label{ga:thm:relative}
The actual cusp boundary map has integral cone homology
\[
H_n(\operatorname{Cone}\kappa_\infty)=
\begin{cases}
0&n=0,1,\\
\mathbb Z&n=2,6,\\
\mathbb Z^2&n=3,5,\\
\mathbb Z^4&n=4,\\
0&\text{otherwise}.
\end{cases}
\]
\end{theorem}
\begin{proof}
In each degree the outgoing differential selects the first $r_{n-1}$
coordinates of the shifted boundary summand.
Its kernel is the full target summand and the remaining shifted coordinates.
The incoming differential identifies its domain coordinates bijectively
with the entire target summand. Therefore the quotient of cycles by
boundaries is freely generated by those remaining shifted coordinates.
Their ranks are $(0,0,1,2,4,2,1)$.
Every elimination uses a unit coefficient, so the quotient has no torsion.
The preceding chain contractions identify this exact integer calculation
with the cone of \eqref{ga:eq:kappa}.
\end{proof}

\section{The first homology of the common attachment}

We specify the remaining real topological data.
Let $\Sigma$ be an oriented sphere with three open disks removed.
Its marked positive meridians satisfy $\gamma_0\gamma_1\gamma_\infty=1$,
with right-to-left path multiplication. The flat torus bundle has monodromies
\[
T_0=\begin{pmatrix}0&-1&-1&0\\1&-1&0&0\\0&0&1&0\\0&0&0&1\end{pmatrix},\qquad
T_1=\begin{pmatrix}-1&1&-1&2\\-1&0&-1&1\\1&1&2&-1\\0&0&0&1\end{pmatrix}.
\]
They satisfy $T_0T_1T_\infty=I$ in the common $e$ basis.

Here are the finite fillings and their markings.
Set $u_0=e_3$, $u_1=-e_1+e_3$, $u_2=-e_1-e_2+e_3$.
Let $X_0$ be their three-torus and let $P_0$ cyclically permute this basis.
Put $U_0=u_0+u_1+u_2=\delta$, where $\delta=-2e_1-e_2+3e_3$.
For the other three-torus $X_1$ use the $f_1,f_2,f_3$ basis and put
\[
P_1=\begin{pmatrix}0&-1&-1\\1&0&0\\0&0&1\end{pmatrix},\quad
U_1=(-1,-1,2)^t,\quad
C=\begin{pmatrix}1&1&0&1\\0&1&0&0\\0&-1&1&0\\0&0&0&1\end{pmatrix}.
\]
The matrix $C$ sends the $f$ marking to the $e$ marking.
For $\rho=e^{2\pi i/3}$ define
\[
\begin{aligned}
N_0&=(X_0\times S^1_t\times D_s)/
 (x,t,s)\sim(P_0^{-1}x+U_0/3,t+1/3,\rho s),\\
N_1&=(X_1\times S^1_t\times D_s)/
 (x,t,s)\sim(P_1x-U_1/4,t-1/4,is).
\end{aligned}
\]
The finite actions are free because their nonidentity powers translate $t$.
The boundary maps are
\[
j_0[x,t,v]=[x-2vU_0/3,t+v/3,e^{2\pi iv/3}],\qquad
j_1[x,t,v]=[x-vU_1/4,t-v/4,e^{2\pi iv/4}].
\]
For $j_1$, first apply $C^{-1}$ to the fibre marking.
Substitution in these quotients verifies the mapping-torus endpoint relations.
Their core retractions, with $[y,\tau+1]=[P_i y,\tau]$, are
\[
[x,t,v]\longmapsto[x+2tU_0,3t+v],\qquad
[x,t,v]\longmapsto[x-tU_1,4t-v].
\]
Let $W$ be the flat bundle with these two fillings attached.
Its remaining boundary is $B_{T_\infty}$.
This defines a compact real manifold without assuming a global complex family.

Choose compatible finite chains $A=C_*(E)$, $L_i=C_*(N_i)$ and
$B_i=C_*(B_{T_i})$. Their common attachment complex is
\[
\mathcal W=\operatorname{Cone}\bigl(
(b_0,b_1)\mapsto(\iota_0b_0+\iota_1b_1,-j_{0*}b_0,-j_{1*}b_1)\bigr).
\]
The cusp boundary map into it is
$\beta(b)=(b_{\infty*}b,0,0;0,0)$.
Use the same boundary chain complex on both arrows of the attachment.
Subdivision and the cellular-to-singular comparisons transport the
cubical formulas and their specified homotopies to this finite diagram.
The actual attachment has the chain model
\begin{equation}\label{ga:eq:global}
\operatorname{Cone}\left((\beta,-\kappa_\infty):
C_*(B_{T_\infty})\to\mathcal W\oplus C_*(N_c)\right).
\end{equation}
The ordinary union comparison has contractible identity-cone kernels.
It therefore identifies \eqref{ga:eq:global} with the integral chains
of the specified smooth union $Y_{\mathrm{top}}=W\cup_{j_\infty}N_c$.
The retraction and \eqref{ga:eq:comparison} replace its second map by $-K$.
For the homotopy in that equation, the final change of coordinates is
$(w,z,b)\mapsto(w,z-Sb,b)$.
Substitution gives a chain isomorphism between the two core cones.

The finite core maps give the fundamental-group relations
\[
\pi_1(W)=\langle\Lambda,\gamma_0,\gamma_1\mid
[\Lambda,\Lambda]=1,\ \gamma_i\lambda\gamma_i^{-1}=T_i\lambda,
\ \gamma_0^3=w_0,\ \gamma_1^4=w_1\rangle,
\]
where $w_0=(4,2,-6,1)^t$ and $w_1=(1,1,-3,-1)^t$.
Indeed the core maps are coverings in the fourth period, of degrees three
and four, and the positive meridians map to the positive and negative core loops.
Their power relations are $w_0=e_4-2\delta$ and $w_1=-Cf_4-\delta$.
Van Kampen gives the displayed presentation.
The images of $T_0-I$ and $T_1-I$ generate
$\langle e_1,e_2,e_3\rangle$ integrally.
Abelianization therefore reduces to $3x_0=c$ and $4x_1=-c$.
The relation matrix
\[
\begin{pmatrix}-1&3&0\\1&0&4\end{pmatrix}
\]
has Smith factors $(1,1)$ and primitive kernel vector $(12,4,-3)^t$.
Thus $H_1(W)=\mathbb Zg$, with $e_4=12g$ and $\gamma_\infty=-g$.

The ordered basis of $H_1(B_{T_\infty})$ is
$(\beta_1,\beta_2,\gamma_\infty)$, where $\beta_1=e_3$ and $\beta_2=e_4$.
For the target basis $(g,\beta_1,\beta_2)$, the actual degree-one
attachment matrix is
\begin{equation}\label{ga:eq:global1}
(b_{\infty*},-k_*)_1=
\begin{pmatrix}0&12&-1\\-1&0&0\\0&-1&0\end{pmatrix}.
\end{equation}
Its determinant is $-1$. It is an isomorphism over $\mathbb Z$.
The degree-zero map is $(1,-1)^t$.
The cone exact sequence now gives
\begin{equation}\label{ga:eq:first}
H_0(Y_{\mathrm{top}};\mathbb Z)=\mathbb Z,
\qquad H_1(Y_{\mathrm{top}};\mathbb Z)=0.
\end{equation}

There is also a precise description of the next required map.
The cusp map is surjective on $H_2$, and its kernel is
$\mathbb Zz_1\oplus\mathbb Zz_2$ from \eqref{ga:eq:kernel2}.
Each class in $H_2(Z)$ has an integral fibre lift, and hence an integral
lift to $H_2(B_T)$. Eliminating that target summand in the cone exact
sequence gives
\begin{equation}\label{ga:eq:next}
H_2(Y_{\mathrm{top}};\mathbb Z)
\cong H_2(W;\mathbb Z)/
\langle b_{\infty*}z_1,b_{\infty*}z_2\rangle.
\end{equation}
The classes on the right are the actual cusp-meridian tori, rather than
two unnamed elements with prescribed ranks.
Their inclusion into the common complex must be evaluated to determine
this group, including any torsion.

For application to a global complex family, a further geometric comparison
must identify its punctured torus bundle with $E$ relative to all three
marked boundaries and the section. At each finite puncture, the lifted
translation on the reduced divisor torsor must extend smoothly, commute
with the cyclic deck action, and have a proper quotient with the displayed periods.
The retained section must realize the conormal comparisons
$e=s^{-2}e'$ at the order-three puncture and $e=s^{-1}e'$ at the order-four puncture.
These are existence and compatibility conditions on maps of the actual
torsors. They do not follow from their integral monodromy matrices.
The local cusp calculation and \eqref{ga:eq:first} leave those conditions
and the remaining homology of the common attachment distinct.

\end{document}
